\chapter{关系基底：高阶结构、体验记录与表示容量}
\label{chp:static_substrate}

上一章定义了具有身份、来源和绑定条件的语义单元。本章进一步讨论这些单元怎样组织成可供更新的关系基底。我们暂时固定结构版本，分别规定对象之间的关系、对象内部的属性空间，以及经历对关联与控制的影响。静态基底并不是停止运行的智能，而是动态分析中用来明确作用对象的一张结构快照：当结构发生增删、合并或重编码，我们需要记录事件，重新检查维度、接口和先前结论的有效域。

我们保留图、单纯复形、局部属性空间和表示迁移等几何工具，但不把它们预先解释为价值作用的自然场。高阶关系必须来自任务语义，体验记录必须保留证据与时间，度量必须说明比较什么。网络结构的优化结果也需要由目标和可行集合推导，不能从一种构成类比断言所有系统必然形成同一拓扑。本章以这些区分为基础，讨论结构构造、经历更新、统计区分度、资源取舍与容量边界，为后续状态传播和结构适应建立可核验的条件。

\section{高阶关系的构造：图、超边与单纯复形}
\label{sec:morpho_elements}

我们从任务相关对象的身份集合出发构造结构。节点可以表示杯子、人员、事件、命题或程序状态，边可以表示邻接、包含、依赖、先后或支持，但这些类型不能混同。一条空间邻接边只说明两个对象接近，不说明它们有因果作用；一条引用边说明某段结论指向某个来源，也不保证来源可靠。世界图的作用是组织系统当前所承诺的关系及其证据，不是把现实世界压缩为一种无类型连通性。不同任务可以共享对象身份，同时使用不同关系层和读出接口。

普通图适合表示成对关系，超图适合把多个对象放入同一个关系实例。例如，搬运任务中人员、工具和货物共同参与某次操作，这个参与关系是三元事件；三个对象两两有关联，并不足以证明它们共同参与同一次操作。我们可以创建事件节点并标记角色，也可以用带角色的超边直接表示。若使用无序超边，需要另存谁实施动作、谁接受动作及谁提供工具，否则关系方向会丢失。图和超图各有表示代价，选择的依据是任务能否保持必要区别，而不是高阶名称天然比低阶名称具有更强解释力。

抽象单纯复形还具有额外约束。设有限节点集合为 $V$，我们定义 $\mathcal K$ 为其中非空有限子集组成的集合族，并要求每个单纯形的所有非空子集也属于 $\mathcal K$。包含一个元素的是零维单纯形，包含两个元素的是一维单纯形，包含三个元素的是二维单纯形。闭包条件保证一个面具有相应的边和节点，却不反过来要求有三条边就一定有一个面。若把图中的所有团自动补成单纯形，我们得到的是团复形构造，它是一项明确的设计选择，不是成对一致自动产生高阶语义的数学结论。

以三人协作签署文件为例，甲与乙能够合作、乙与丙能够合作、甲与丙也能够合作，并不保证三人能够共同满足同一签署规则。权限、时间和资源冲突可能只在三方同时参加时出现。若我们给这个三角形填面，就等于新增一种高阶兼容承诺，必须由三方约束的检查结果支持。同样，三个命题两两可兼容也不能保证联合可满足。我们把高阶单元视为经过证据或规则确认的联合关系，拒绝以低目标值或图形闭合代替语义验证。缺少联合证据时，保留未填面的边界是合理状态，不是必须消除的结构缺陷。

为了在复形上计算，我们为每个单纯形选择一个方向，定义相应的链空间和边界矩阵。链空间中的系数是无量纲数值，方向只是计算约定，不自动表示现实因果方向。设 $B_k$ 是从第 $k$ 阶链到第 $k-1$ 阶链的边界矩阵，其尺寸为 $n_{k-1}\times n_k$。按交替符号删除顶点的定义，相同低阶面会在连续两次取边界时出现两次且符号相反，因此 $B_{k-1}B_k=0$。在最低阶或最高阶没有相应单纯形时，我们把相应边界矩阵视为零映射。这个结果依赖复形闭包和一致方向；缺少面或错误方向的实现需要先修正数据结构，不能继续使用标准边界恒等式。

在选定欧氏内积后，我们可定义
\[
 L_k=B_k^{\mathsf T}B_k+B_{k+1}B_{k+1}^{\mathsf T}.
\]
对任意 $x\in\mathbb R^{n_k}$，有 $x^{\mathsf T}L_kx=\|B_kx\|^2+\|B_{k+1}^{\mathsf T}x\|^2\geq0$，因此这个算子对称半正定。我们可以用它约束某类数值活动的平滑或一致性，但不能因此宣布高阶关系为真。非欧氏加权内积需要相应重建伴随算子，带类型或有向关系也不能直接塞入这一对称公式。计算算子和语义约束是相互连接的两个层次，前者的数值性质不替代后者的任务含义。

有限复形并不因为节点数量增加就必然逼近黎曼流形。若要采用连续极限，我们至少需要指定采样空间、邻域构造、尺度变化、权重及收敛目标，并检查边界和奇异结构。某些关系图本身没有自然的连续几何对象，仍然可以直接使用离散算子完成任务。我们因此优先给出有限结构上的定义，再在条件充分时引入连续近似。这样的顺序既保留几何工具的价值，也避免把未经建立的极限当成后续动力学的默认事实。

\section{世界结构与体验记录：关联、控制和路径依赖}
\label{sec:section_8660}

我们把世界结构与体验记录区分为两种数据责任。世界结构保存当前对象、关系及相应来源；体验记录保存系统经历的观测、动作、预期、结果、评价和修订。体验可以通过更新规则改变结构权重、检索次序或控制偏好，但它并不因被称为体验图就天然成为一种联络。若经历采用图组织，节点可以是事件，边可以是时间先后、同一对象延续或证据依赖；若采用日志组织，也可以提供同样的更新信息。它们的共同要求是保留时间、主体和来源，而不是必须服从一种指定的几何本体。

饥饿状态下苹果更值得优先获取，饱腹后这种优先性可能下降。我们可以把这一变化放入目标权重、动作效用或注意门控中，并同时保留苹果的身份与可测属性。价值变化不要求对象向量旋转，也不要求所有逻辑路径的长度改变。若任务读出采用向量编码，价值可以成为一个字段，但字段与控制器之间的关系需要另行规定。把世界事实和主体评价分开，还允许多个主体共享对象记录而具有不同选择；否则一方的偏好可能被误写为所有主体必须接受的环境事实。

若系统需要跨局部表示接口传递属性，我们可以定义边上的线性映射。设节点 $i$ 和 $j$ 的属性空间分别为 $F_i=\mathbb R^{d_i}$ 和 $F_j=\mathbb R^{d_j}$，映射 $T_{ji}:F_i\to F_j$ 为 $d_j\times d_i$ 矩阵。该映射可以转换单位归一化后的特征、抽取共同字段或生成接收节点所需的消息。不同维度时通常不可逆，故我们称其为接口映射，而不自动把它当作规范坐标变换。语义迁移还需要身份、时间及读出约束，矩阵乘法尺寸正确只是最低条件。

在共同固定维数的特殊情形，路径 $i_0\to i_1\to\cdots\to i_m$ 的接口合成为
\[
 T_{\gamma}=T_{i_m i_{m-1}}\cdots T_{i_1 i_0}.
\]
乘法顺序由消息经过节点的先后确定。若闭路径回到原节点而 $T_{\gamma}\ne I$，同一属性向量经过该接口序列后可能改变。这是可定义的路径依赖，不说明改变一定是成长或创伤。若所有映射仅来自可逆坐标重标记，即 $T_{ji}=R_jR_i^{-1}$，逐项抵消得到任意闭路径 $T_{\gamma}=I$。因此回路中出现变化意味着存在额外更新、信息损失或不一致映射，需要检验其来源，而不能只凭闭路径的形状解释。

有些路径依赖来自事件记忆的积累，而非接口矩阵。系统完成访问甲地点、访问乙地点、返回甲地点的行程后，空间位置回到起点，经历记录却增加了事件。若内部状态同时包含位置和日志，整体状态没有回到初值，这是记录定义的直接结果，机器程序也可以如此。反过来，有学习能力的系统也可能在某些可逆处理之后恢复原状态。我们不以循环后是否改变来划分机器与智能，而考察改变是否保存有用信息、是否遵守任务约束以及是否能影响后续适应。

体验对结构的更新还必须区分单次结果与稳定规律。一条路径执行失败可能来自工具故障、观测延迟或偶然外部扰动，不能立即把全部相关边标成不可通行。我们可以保存事件并更新局部置信，等待重复证据或进行干预后再修订结构模型。涉及价值判断时，评价来源与任务承诺同样需要记录；降低内部不适不能自动成为修改事实或取消安全约束的理由。错误经历也可以长期积累，故路径依赖的存在并不是学习正确性的证明。

在 AgentOS 中，体验事件主要进入情景记忆 $E$，与当前 $h(t)$ 的对象记录保持身份和版本关联；可复用规则进入 $\Theta$ 时需要明确抽象条件。自我相关状态 $\Psi$ 可以调节优先级，SPC和TCE关注运行状态及控制，但不因此替代事实来源。外部记忆可以保存更完整的历史，使工作界面只加载当前任务所需的片段。我们把这些接口看作一个工程实例：一般HSF-HD模型只要求声明状态、结构、来源和更新机制，不要求全部智能都使用相同的记忆组件名称。

\section{局部属性空间与统计区分度的条件}
\label{sec:section_7949}

每个节点所携带的属性应先按类型组织，再决定是否采用向量空间。颜色、重量、温度、情绪效价和文本编码并非天然同类量；某些字段是连续测量，某些是离散类别，某些是多值集合或缺失状态。我们可以在模型中把连续字段按参考尺度归一化，组成有限维向量，也可以把类别编码与数值字段拼接。但拼接只是接口形式，不会消除原字段的单位与来源。节点类型不同，属性维度和合法集合也可以不同，这时整个状态空间更适合按类型分区，而不是强制成为同一个固定纤维空间。

在一个明确的实内积空间中，正交基可以方便展开向量。若 $q=\sum_a q_a e_a$，各基向量的内积为零，只说明所选坐标的几何关系；它不说明颜色和重量的随机变化统计独立，也不说明感知通道彼此不影响。属性之间的依赖可以由协方差、条件分布或结构关系记录，正交化可能改变字段解释。活动门控也应与属性数值分开：红色字段的数值可以稳定，当前注意程度则改变。把活动系数当作全部内容会让低活动记录看似丢失事实，即使它仍可从记忆中准确读出。

统计区分度提供另一种几何工具。我们先定义一个有限结果集合和参数域，参数 $\theta\in\mathbb R^p$ 按参考尺度归一化。假设每个结果的概率 $p_\theta(x)$ 为正，随参数可微，支持集固定。评分向量为 $s_\theta(x)=\nabla_\theta\log p_\theta(x)$，Fisher矩阵定义为 $I(\theta)=\sum_xp_\theta(x)s_\theta(x)s_\theta(x)^{\mathsf T}$。其尺寸为 $p\times p$。对任意方向 $v$，有 $v^{\mathsf T}Iv=\sum_xp_\theta(x)(v^{\mathsf T}s_\theta(x))^2\geq0$，故半正定性直接来自平方和；正定性则要求每个非零方向都能改变某些结果的对数概率。

如果两个参数只通过它们的和影响概率，改变其中一个并反向改变另一个就不改变分布，Fisher矩阵在该方向退化。这说明度量退化可能来自参数不可辨识，不能把它解释成概念没有内容。反过来，一个属性向量具有很大样本方差，也不意味着所有参数方向都可区分。分布之间的全局接近与某一点的局部Fisher值同样不是同一问题：两个相距很远的参数可以产生相同分布，而局部评分仍然不为零。我们因此不由某个散度较小就推出Fisher矩阵趋于零，更不把这类统计量直接解释为语义体积或价值曲率。

在有限支持且具有固定基准函数的指数族中，我们可以进一步推导协方差形式。设充分统计向量为 $\phi(x)\in\mathbb R^p$，模型定义为
\[
 p_\theta(x)=\exp\bigl(\theta^{\mathsf T}\phi(x)-A(\theta)\bigr)b(x),\qquad
 A(\theta)=\log\sum_x b(x)e^{\theta^{\mathsf T}\phi(x)}.
\]
这里 $b(x)>0$，所有指数项无量纲，有限求和允许逐项求导。由归一化项求导得到 $\nabla A=\mathbb E_\theta[\phi]$，所以评分为 $\phi-\mathbb E_\theta[\phi]$，再代入定义得到
\[
 I(\theta)=\operatorname{Cov}_\theta(\phi).
\]
这是一类声明了概率模型的结果，不依赖热力学类比。若额外引入缩放参数，协方差前的系数必须由评分中的缩放推导；任意输入模型都不能无条件套用同一协方差公式。

在更高阶可微条件下，局部KL散度的二阶展开为 $\frac12\delta^{\mathsf T}I(\theta)\delta$ 加高阶余项。我们先对第二个参数求导，利用评分均值为零使一阶项消失，再由归一化恒等式把二阶项写成Fisher矩阵。这个近似只适用于参数邻域和声明的支持条件；若概率为零、支持改变或不可微，需要采用其他分析。使用Fisher作为更新几何，还要处理退化方向与正则化，不能仅凭名称就得到可逆度量或稳定控制。

因此，学习到的向量空间、任务定义的距离和统计模型的Fisher几何是三个可以联合但不必相同的对象。词向量可能捕获任务有用的共现模式，具体编码是否保持某种统计距离，需要明确映射并验证误差，而不是直接宣布全部嵌入都是Fisher映射。我们在本章用统计几何解释参数可区分性，用有类型字段解释属性，用接口映射解释迁移。这样的分工使模型能够选择合适工具，并保留工具失效时可定位的原因。

\section{网络结构的资源取舍与受限编辑}
\label{sec:section_6753}

我们研究模块、短路径和长程接口，是因为它们影响信息访问、验证与维护，而不是因为一种通用构成目标必然把所有网络推向同一形态。连通程度较高可以降低某些路径长度，却会增加存储、通信、权限检查和错误传播的负担。局部模块便于封装知识与限制更新范围，却可能使跨模块任务需要额外接口。是否值得新增长程连接，取决于任务分布、访问频率和维护代价；同样的网络在一种工作负载下有效，在另一种工作负载下可能浪费资源。

我们先声明固定节点集合和允许的边集合，令 $G$ 表示其中一个可行结构。对指定的访问请求分布，我们定义平均访问损失 $D(G)$；对每条实际边定义存储或维护成本，再得到总成本 $C(G)$；对可能造成越权传播、身份混淆或失效扩散的连接定义风险指标 $R(G)$。不同指标经参考尺度归一化后组合为 $J(G)=\alpha D(G)+\beta C(G)+\gamma R(G)$，权重非负且由工程需求确定。不可达请求需要明确的有限惩罚或作为硬约束排除，不能在计算时一面使用无穷距离一面宣称目标有限。

一个简单反例说明目标并不自动产生模块化小世界。若节点间所有边成本相同、允许完全图，并且只优化连通图的平均跳数，完全图使每对不同节点只需一步，因而达到该目标的最小可能值。若只优化正边成本且不要求连通，空图达到最小值。若只优化边成本但要求连通，所有具有相同边成本的生成树都可能最优，它们的平均路径长度和聚类程度却不同。三个目标使用同一组节点而产生不同结构，足以说明结构结论必须由具体取舍和约束支持。

我们可以把模块化和长程接口纳入候选集，也可以对任务群组指定局部访问偏好，再评估它们是否降低综合目标。若某个接口连接两类频繁协作的模块，它可能降低访问损失；若接口被用于不可信信息的快速扩散，风险项可能增加。枢纽能缩短某些路径，却也可能形成吞吐瓶颈与单点故障。因此新增边的净效果应通过增量评价或重计算确认，不能仅凭“捷径”名称判断有利。图中的边还需保留类型，不能把所有关系都当作可无条件传递消息的通信线路。

在有限候选结构集合上，我们可以定义严格下降的编辑程序。每次只接受满足约束且使目标下降至少 $\delta>0$ 的编辑；假设目标下界为 $J_{\min}$，则从初值出发，接受的编辑次数至多为 $(J(G_0)-J_{\min})/\delta$ 的整数下取整。推导只需要累计下降与下界，但终止后得到的只是相对于允许编辑集合的局部不可改进结构。它不保证全局最优，也不保证具有某种预定拓扑。若目标、任务分布或节点集合随时间改变，这一固定目标终止论证需要重新建立，不能把历次分属不同目标的下降相加。

结构编辑和活动更新还应分别提交。新增一条关系需要来源、权限、版本和后续失效处理，活动传播只改变固定结构上的当前数值。删除一个节点后，关联边、属性绑定、索引、缓存和正在执行的计划必须同步检查；否则看似简洁的网络可能包含悬空引用。我们可以在事务快照上验证候选编辑，提交后通知依赖项重新读取。内部回滚可以恢复记录版本，但不会撤销已经发生的现实动作，因此现实执行应通过独立接口和反馈确认。

以知识库维护为例，两个领域都使用同一术语，却具有不同定义。创建跨域边可以提高检索速度，也可能把一种定义误用于另一种任务。我们可以采用带类型的转换接口，保留源域、目标域及适用条件；需要时再根据使用频率决定缓存或预计算。长期低使用率的边可能被压缩或卸载到外部记忆，仍然保留可恢复来源。这是资源与任务的持续取舍，不是网络趋向低能状态的自然必然。AgentOS的CCM和Offloading可以实例化此类管理，但一般理论不要求相同的组件安排。

验证时我们分别测量路径与延迟、结构存储、接口维护、错误传播和任务完成率。较短图路径不一定产生较低实际延迟，因为跨设备通信和验证可能更昂贵；较高聚类也不一定意味着更一致的推理，因为重复来源可能造成伪证据增强。我们把模块化与捷径保留为可检验的候选设计，要求通过工作负载和失效场景评估它们，而不把一个结构统计指标当作智能水平或普遍稳定性的代理。

\section{拓扑诊断、有效维数与能力边界}
\label{sec:section_4716}
\label{sec:section_5497}

拓扑工具可以描述结构中哪些关系闭合、哪些缺少高阶填充，但不能独立定义意义、自我或反思。我们在选定系数域的有限复形上定义第 $k$ 阶同调为边界算子核对高一阶边界像的商空间，贝蒂数是这个商空间的维数。因为连续取边界为零，高一阶边界像包含于核，商空间才有定义。对实数系数的有限链空间，可用矩阵秩写成
\[
 \beta_k=n_k-\operatorname{rank}B_k-\operatorname{rank}B_{k+1}.
\]
这个公式来自核维数的秩—零化度关系以及边界像维数，不需要把环解释为反思机制。数据结构、填充规则和系数域改变时，贝蒂数也可能改变，报告结果时应同时报告这些选择。

三角形的例子提供直接检验。只有三条边和三个节点而未填面时，边界环形成一个一阶同调类；加入同一三角形的面后，该环成为面的边界，一阶贝蒂数减少。任务中的甲指向乙、乙指向丙、丙指向甲，则是一个有向依赖循环，其含义取决于边类型和执行规则。若它表示带延迟的反馈，可能支持稳定调节；若它表示互相等待的资源锁，可能形成死锁。把它转为无向复形后得到的洞既不区分这两种机制，也不保留全部方向信息。拓扑结构与动态反馈必须联合分析，不能互相替代。

缺少面可以表示尚未确认的高阶兼容，也可能只是采样不足或建模约定。我们可以通过补充观测检查缺失关系，或在一组尺度上比较结构摘要的稳定性，但稳定摘要仍不自动具有某个心理含义。自我相关状态需要指称、目标、历史和控制接口，不能仅因存在二维腔就认定其中居住着自我；意义需要任务读出和使用过程，不能把加权贝蒂数之和当作已证明的意义总量。若将这些结构统计量作为经验指标，应先指定预测目标、数据与反例，而不是把指标名称升级为守恒定理。

动态吸引子也需要独立定义。一个有界状态在指定更新下趋近某个集合，属于动力学性质；节点被标记为高价值，则是目标或偏好字段。高价值对象不一定吸引实际状态，因为权限、资源或其他承诺可以阻止访问。我们不使用趋于正负无穷的价值来解释所有生存或死亡相关选择，而允许有限偏好与硬约束共同工作。目标极小点是否稳定、是否可达和是否有益，分别需要分析更新、初值与环境结果。拓扑洞的存在不会自动证明其中某个状态是吸引子。

维度分析首先区分关系的最大阶数、嵌入向量维数、属性字段数量、统计有效秩和任务所需信息量。它们可以相关，却不是同一个数。用很多坐标嵌入一个简单图不会新增真实关系，增加字段也可能只是复制已有信息。反过来，复杂任务可以通过程序、外部记忆和分层组合处理，而不要求一个固定界面同时容纳全部状态。我们因此不由某个几何维数直接推断智商、情感体验或抽象能力，也不设定十一维为普遍最优值。具体能力应由任务和有效读出决定。

有限精度给出一项清楚的容量上界。若一个内部编码有 $d$ 个坐标，每个坐标只取至多 $2^b$ 个可区分值，且没有额外存储，那么编码总数至多为 $2^{bd}$。要对 $N$ 个任务等价类进行无误、单值区分，必须满足 $N\leq2^{bd}$，否则由抽屉原理至少两个类别共享编码。该界依赖有限精度和无额外通道的条件，只是必要容量条件，不保证学习算法能够使用全部编码，更不保证编码具有正确语义。允许外部记忆、时间序列或交互式查询后，必须把这些通道一同计入，而不能继续沿用孤立编码的界。

对线性压缩也有对应边界。设任务相关变化包含一个 $r$ 维线性子空间，压缩矩阵把它映射到 $d$ 维空间且 $d<r$。限制到该子空间的映射必有非零核，因此存在两个不同输入得到同一压缩结果；若任务需要区分它们，就无法从压缩结果唯一恢复。若任务本来不区分这些差异，压缩却可以合法完成。由此决定是否需要增加维度的不是输入全部自由度，而是任务必须保留的区别、允许误差和可使用接口。扩大表示能缓解某些瓶颈，也会增加训练、检索及控制成本，收益没有无条件单调性。

例如，同样的红色词汇表可以服务粗略报警或细致色彩管理，所需分辨率不同；混合情绪记录可以包含多个评价分量，但向量维数不能单独证明主观体验。硬件的三维堆叠、总线宽度和计算布置影响通信及吞吐，需要按实际实现测量，不能把芯片空间维数等同于抽象关系维数。我们在工程上先明确任务区别和接口预算，再选择关系阶数、属性精度、外部访问与并行结构。这条路线使表示容量分析保持严谨，也为后续智能状态动力学提供真实的受限空间，而不靠维度意象许诺通用能力。
